% ---------------------------------------------------------------------------
% Draft subsection for Report_metrics.tex.
% Place it immediately after \subsection{Regularity of the solution and the
% left--right bifurcation} (sec:bif), of which it is the continuation.
%
% Numbers come from viz/out/homotopy_flips.json and viz/out/homotopy_lock*.json
% through the macros in homotopy_macros.tex and the tables tab/homotopy.tex and
% tab/homotopydelta.tex.  Nothing below is typed by hand.
% ---------------------------------------------------------------------------

\subsection{Where the discrete choice is actually made}
\label{sec:homotopy}

Section~\ref{sec:bif} closes on a negative result: at the deployed weight the cost landscape does not
bifurcate, the minimiser is unique, and the control law is regular in $x_0$. That settles what the program
does, and it raises the question the program cannot answer. Passing an obstacle on the left and passing it on
the right are distinct manoeuvres; something in the stack decides between them. If the nonlinear program has
a single minimum, the decision is not taken there.

The course anticipates the situation. Discussing mixed-integer programs, the notes observe that a large
number of integer variables is to be avoided and that, in optimal control, one can instead
\emph{pre-select a switching strategy with continuous tuning parameters}, trading degrees of freedom against
the complexity of the resulting program.\footnote{Lecture notes, \S4.2.6.} That is precisely the architecture
of this stack, and it has never been claimed as a design choice: ``which side do I pass on'' is the integer
variable the formulation refuses to carry, and it is delegated to the \astar{} layer of
Section~\ref{sec:astar}, which re-plans on a freshly inflated grid every few control cycles. The price of
keeping the program smooth is that the discrete decision is made somewhere that has no obligation to be
consistent with itself.

\paragraph{Measuring the consistency of the delegated choice.}
Obstacle returns are grouped into connected clusters, and for every cluster the reference passes, we record
which side it passes on. The side is read as a lateral offset at the cluster's longitudinal station, in the
frame fixed by the start pose and the goal; the landmark is the obstacle point nearest the reference, which
behaves alike on a compact pillar and on a wall running parallel to the direction of travel. A \emph{flip} is
a cluster engaged by two consecutive references whose side differs between them; clusters engaged by only one
of the two are ignored, since entering or leaving the planning window is not a change of mind.

Two more obvious definitions were tried first and both proved to measure the wrong thing. The sign of the
cross product between the path tangent and the obstacle at closest approach changes whenever the path merely
curves. The winding angle about the obstacle --- the homotopy invariant of a \emph{fixed} path --- changes as
the robot advances past it, so under a receding horizon it encodes longitudinal progress rather than a
decision, and made every cluster of one scenario flip on the same cycle. Reading the offset at a fixed
station is insensitive to both.

\restab{homotopy}

On the trap scenario the reference changes side $\resHomoFlipsUtrap$ times in $\resHomoReplansUtrap$
re-plans, and the recorded offset is $\resHomoOffsetUtrap$~m at every one of them: the sign alternates while
the magnitude does not, so each flip translates the reference by twice that distance across a stationary
obstacle. On the corridor the count is $\resHomoFlipsCorridor$, and the offsets show why --- the planner is
routing around the \emph{outside} of the corridor walls and alternating which side it leaves by. The two
control scenarios measure zero, correctly: through a narrow gap and around a single pillar the choice, once
made, is never revisited.

\paragraph{The consequence is gated by the horizon.}
At the deployed tuning these flips cost nothing. With the layer of the next paragraph enabled, the executed
trajectory of the trap scenario is identical to the baseline in every digit --- same length, same clearance,
same number of cycles --- while the flip count falls to zero. The explanation is the horizon of
Section~\ref{sec:horizon}: as already noted in Section~\ref{sec:fhocp}, $N\Delta t=\resHorizonSeconds$~s at a
cruise speed of $\resVref$~m/s spans well under a metre of travel, and the two routes do not separate within
that distance. The reference moves several
metres sideways and the controller never sees it, because the divergence lies beyond the last predicted node.
The instability is real and, at this tuning, invisible --- which is a different statement from harmless, and
the next paragraph separates the two.

\paragraph{Pre-selection with hysteresis.}
The remedy follows the notes' prescription literally: make the switching strategy explicit and give it a
continuous parameter. At each re-plan the new \astar{} route is compared with the committed one. If the
signatures agree, nothing further is computed. If they disagree, a route in the committed class is recovered
--- by re-anchoring the route already being followed at the current pose, or, if that route is no longer
usable, by re-planning behind a temporary barrier on the side the challenger chose --- and both are scored
against the real obstacles with
%
\begin{equation}
  J_{\text{route}} \;=\; \int \mathrm{d}s \;+\; w_{\text{clear}} \int \max\!\big(0,\, d_{\text{ref}} -
  d_{\text{obs}}(s)\big)\,\mathrm{d}s ,
  \label{eq:routescore}
\end{equation}
%
the first term being \astar{}'s own path cost and the second the gap-width criterion, written so that
squeezing past an obstacle is paid for in the same units as a detour. The challenger is adopted only if it
beats the incumbent by more than $\delta$. Nothing in the nonlinear program changes: the layer is upstream of
it, the reference handed to the controller has the same form it always had, and $J^\star$ on any single cycle
is untouched.

\restab{homotopydelta}

Three readings of the sweep. A margin of exactly zero is not enough --- the challenger then wins on any
improvement however small, and the flips survive; the hysteresis has to be strictly positive to exist at all,
which is the content of the notes' ``continuous tuning parameter''. From $\delta=\resHomoDeltaSat$ upward the
behaviour is constant, so the deployed value sits on a plateau rather than at a fitted optimum. And with a
prediction horizon long enough for the controller to see the divergence
($N=\resHomoNlong$), removing the flips is worth $\resHomoPathOff \to \resHomoPathLock$~m of executed path,
$\resHomoClearOff \to \resHomoClearLock$~m of minimum clearance, and $\resHomoSuccOff \to
\resHomoSuccLock$~\% of solved cycles.

That last row is the substantive result, and it inverts the usual reading of Section~\ref{sec:horizon}. There
the short horizon was shown to be dominated on time and clearance; here it is what conceals a defect in the
layer above. A controller that cannot see far enough to be affected by an inconsistent reference is not
thereby robust to it.

\paragraph{What is not established.}
The trade-off the notes describe --- degrees of freedom against complexity, here holding a class against
refusing a better route --- was not observed. Raising $\delta$ to $\resHomoDeltaMax$, far beyond any route
cost in these scenarios, degraded nothing, because in this set the committed class is never the worse one.
Demonstrating the cost of hysteresis requires a scenario in which the world changes after the commitment is
made, and that experiment has not been run. On the corridor a single flip survives the layer, so the
mechanism is effective rather than complete. Finally, all of the above is measured in the offline harness,
which tracks a look-ahead setpoint with a proportional controller; as in Section~\ref{sec:robust}, fine
differences between references are attenuated by that outer loop, and confirmation on a recorded mission
remains to be done.
